\section{A total-function completion obstruction and AND--OR comparison}
\label{app:andor-transfer}
These supporting results fix the physical queries and unit prices. We
prove an exact minimum over every compatible total Boolean function on
one complete source, then examine a specific transfer of Ito--Suzuki's
total-tree result~\cite{itosuzuki2025collapse}. Other reductions and
changes of menu or cost remain outside these claims.

\subsection{A complete source and a compatible total function}
Use eight records and the four indexed queries
$q_i=(2i,2i+1)$, $i=0,1,2,3$. Under label zero impose the SAME partition
$(0,1,0,1,0,2,0,3)$ and $\DIFF(5,7)$. Under label one impose the SAME
partition $(0,0,0,1,0,1,0,1)$ and no DIFFERENT constraint.
Retain every satisfying equivalence partition. Encode a word by
$w=\sum_{i=0}^{3}2^i x_i$. The full answer supports are
\[
 W_0=\{0,3,4,7,8,11\},\qquad W_1=\{1,15\}.
\]
Indeed, the first prior gives $x_0=x_1$ and forbids $x_2=x_3=1$;
the second gives $x_0=1$ and $x_1=x_2=x_3$. Assigning each remaining
block a distinct fresh class proves sufficiency. This is the existing complete eight-input source in the four-query proof.

The total function
\[
 g(x)=(x_0\wedge\neg x_1)\vee(x_2\wedge x_3)
\]
is compatible. In variables $(x_0,1-x_1,x_2,x_3)$ this is a positive AND--OR tree.
Negation decodes the same paid query; it is not a physical symmetry in
the source catalogue.

\subsection{The forced hard-product distribution leaves the promise}
Let $S$ be uniform on the eight original words, and let $L,R$ be the two
AND child values of $g$. Their joint law is
\[
 \delta(0,0)=\frac34,\quad
 \delta(1,0)=\delta(0,1)=\frac18,\quad\delta(1,1)=0.
\]
For a unit-cost AND of two literals, a distribution conditional on
output zero assigns probabilities $p,q$ to the two single-zero assignments
and $r$ to the double-zero assignment. The two directional orders cost
$1+q$ and $1+p$, so
\[
 \min(1+p,1+q)\le 1+\frac{p+q}{2}
      =\frac32-\frac r2\le\frac32.
\]
The unique least-favourable distribution has $r=0$, $p=q=1/2$;
conditional on output one there is only one assignment. Thus the hard
child laws required in the cited Lemma~4.3 are forced in this example:
the left zero law is uniform on actual bits $00,11$, and the right
zero law is uniform on $10,01$. Their one laws are $10$ and $11$,
respectively.

Combining these laws independently conditional on the child values,
as in the cited Definition~4.1, gives the following positive masses:
\[
 S^\delta(w)=
 \begin{cases}
  3/16,&w\in\{4,7,8,11\},\\
  1/16,&w\in\{5,9,12,15\}.
 \end{cases}
\]
Words $5,9,12\notin W_0\cup W_1$ carry exactly $3/16$ mass
outside the original promise. Uniqueness forces this escape for the
specified completion, $\delta$ and hard-child requirement. The product
remains legal on the cited paper's full cube; its preservation of this
relational promise would be an extra, false premise.

\subsection{The values of this promise and this completion differ}
At unit prices, the existing four-query certificate specializes to
$R_{\mathcal A}=R_{\mathcal P}=5/2$ on this source. This value also has
a short direct proof. Opposite-label single-bit neighbours force queries
$0,1$ on word $1$ and $2,3$ on word $15$. The first query of any correct
tree is therefore outside the mandatory pair of one of these words.
Its cost is at least three there and at least two on the other word.
The prior uniform on $1,15$ gives lower $5/2$, including for mixtures.
Zero-prior inputs remain required for correctness.

For the upper bound let $N_i(L,R)$ query $i$ and take $L,R$ on answers
zero and one, respectively. Mix the two legal priority trees
\[
 \begin{split}
 T_1&=N_1\bigl(N_0(0,1),N_2(0,N_3(0,1))\bigr),\\
 T_2&=N_2\bigl(N_1(N_0(0,1),0),N_3(0,1)\bigr)
 \end{split}
\]
equally, with orders $(1,0,2,3)$ and $(2,1,0,3)$. On words
$(0,1,3,4,7,8,11,15)$ their costs are $(2,2,2,2,3,2,2,3)$ and
$(3,3,2,2,2,3,2,2)$, so every expected cost is at most $5/2$.
These are the existing policies $P_5,P_{14}$ on a certified tie face;
no new minimax solve is used.

For the specified total $g$, both full-cube values equal $3$.
For the lower bound give words $4,7,8,11$ equal probability. Independently
in each child, the unique zero literal has a uniformly random position.
Full-cube correctness requires finding a zero in each AND. The first
query in either child costs one, and with probability $1/2$ forces its
second query, even after observations in the other independent child.
Hence every adaptive algorithm has expected cost at least $3$.
The eight directional seeds uniformly order the root's children and each
child's leaves. When both children are false, each costs at most $3/2$
in expectation; when both are true, cost is $2$. When exactly one is true,
the uniform root order gives at most $(2+(3/2+2))/2=11/4$. Thus the
full-cube priority upper bound is $3$.
Restriction gives a valid but non-sharp promise upper bound. The next
subsection proves that the full-cube minimum remains $3$ over every
compatible total function on this same menu, including nonunate ones.

\subsection{The minimum over all compatible total functions}
\label{sec:unit-common-completion}
Let $W=W_0\cup W_1$ be this eight-word promise. Write $\Comp(W)$ for
all compatible total Boolean functions on the same four coordinates.
For $f\in\Comp(W)$, $R_{\mathcal C}(f;\mathbf1)$ requires correctness
on the full cube; every seed computes the same $f$. The notation
$R_{\mathcal C}(W;\mathbf1)$ charges only the original promise.

\begin{proposition}[Exact unit-cost common-completion minimum]
\label{prop:unit-common-completion}
On this complete source, with unit query prices and zero-error
input-independent private randomness,
\[
 \begin{aligned}
 \min_{f\in\Comp(W)}R_{\mathcal A}(f;\mathbf1)
 &=\min_{f\in\Comp(W)}R_{\mathcal P}(f;\mathbf1)=3,\\
 R_{\mathcal A}(W;\mathbf1)&=R_{\mathcal P}(W;\mathbf1)=\frac52.
 \end{aligned}
\]
The minima include nonunate functions. Thus the common-total-completion
minimum is exactly $6/5$ times the promise value on this fixed source.
\end{proposition}
\begin{proof}
The promise value and the full-cube upper bound $3$ for the displayed
$g$ were proved above. We prove that every compatible total $f$ has
$R_{\mathcal A}(f;\mathbf1)\ge3$.

Let $w^{(q)}$ denote word $w$ with coordinate $q$ toggled, and define
$s_f(w)=|\{q:f(w)\ne f(w^{(q)})\}|$. If these labels differ, any
deterministic tree computing $f$ must query $q$ on $w$: otherwise both
inputs have the same purchased transcript and output. Therefore every
zero-error seed costs at least $s_f(w)$ on that same input, and
\[
 R_{\mathcal A}(f;\mathbf1)\ge\max_w s_f(w).
\]
For an arbitrary seed probability space, discard the finite union of
erroneous null seed sets over all $16$ inputs first. This preserves
expected costs. The sensitivity lower bound is standard decision-tree
reasoning; no new general lower-bound method is proposed.

Suppose all sensitivities are at most two. The eight original labels
then force every other value by the following six implications.
\begin{enumerate}
 \item Word $1$ has value one and its neighbours $0,3$ have value zero.
 Its other neighbours $5,9$ must have value one.
 \item Word $15$ has value one and neighbours $7,11$ have value zero.
 Its other neighbours $13,14$ must have value one.
 \item Word $7$ has value zero and neighbours $5,15$ now have value one.
 Its neighbour $6$ must have value zero.
 \item Word $11$ has value zero and neighbours $9,15$ now have value one.
 Its neighbour $10$ must have value zero.
 \item The four neighbours $0,3,6,10$ of word $2$ all have value zero.
 Value one at $2$ would give sensitivity four, so $f(2)=0$.
 \item Word $14$ now has value one and neighbours $6,10$ have value zero.
 Its neighbour $12$ must have value one.
\end{enumerate}
All $16$ values are fixed. Their one-set is
$\{1,5,9,12,13,14,15\}$, exactly the one-set of $g$. Thus every
compatible $f\ne g$ has an input with at least three mandatory queries
and hence randomized value at least three. For $g$ itself, the preceding
independent-child hard prior gives lower three even though its maximum
sensitivity is only two. The legal priority upper three for $g$ and
$R_{\mathcal A}\le R_{\mathcal P}$ complete both minima.
\end{proof}

This excludes a value-preserving common-total-function route on the
unchanged menu at these prices. Promise-correct seed trees may have
different off-promise completions; their value $5/2$ is therefore
consistent with the common-completion minimum $3$. The off-promise
words used in the proof are correctness obligations for $f$, not added
original source inputs. The result gives no unequal-price statement,
general $6/5$ bound, or change to the five-to-seven randomized additive
adaptive/priority interval. Generic completion phenomena and methods
remain established context; exact literature priority is not certified.

\subsection{Depth-first and global-order sets are incomparable}
For $h=(x_0\wedge x_1)\vee(x_2\wedge x_3)$, global order $(0,2,1,3)$
is a legal priority policy with short-circuit skips. Its prefix $(0,2,1)$
on the all-one input interleaves unfinished children, violating depth-firstness.

Conversely, query $0$ first. If it is zero, evaluate the other child in
order $(2,3)$. If it is one, query $1$; on a zero answer evaluate the
other child in order $(3,2)$, and otherwise stop with one. Short-circuit
each AND. This policy is depth-first and correct on all $16$ words.
At words $0,4$ its paid masks are $5,13$: purchasing $3$ depends on the
answer at $2$, and these words differ only at $2$. Any causal priority
policy with this entire paid-set vector must put $2$ before $3$.
At words $1,9$ the masks are $11,15$, forcing $3$ before $2$ by the same
argument. No global order has this vector. This statement concerns a
specific policy vector; it does not assert unequal class optima for $h$.

Incomparability does not contradict the valid unit-cost total-domain chain
\[
 R_{\mathcal A}\le R_{\mathcal P}\le R_{\rm dir}=R_{\rm DF},
\]
because the last comparison uses the smaller directional embedding and
the cited value theorem, rather than a depth-first/priority inclusion.
A sufficient transfer route to a promise $W$ would supply a compatible
total tree $T$ on the same menu/prices, a valid directional/depth-first
equality at those prices, and $R_{\rm DF}(T)\le R_{\mathcal A}(W)$.
The resulting sandwich proves equality on $W$. On our unit-cost source,
Proposition~\ref{prop:unit-common-completion} rules out that sharp route
through any compatible total tree on the unchanged menu. Other
reductions and changes of the information or cost contract remain open.

\paragraph*{Verification scope.}
Two implementations reconstruct the complete source and check all
$16$ product masses, the policy witnesses and the precedence cycle.
The analytic arguments establish the hard-child uniqueness and lower
bounds. The six fixed sensitivity implications were additionally checked
without enumerating or optimizing completions. Their universal-policy
lower bound is the written proof, rather than a conclusion from fixtures.
These checks establish neither human review, formal certification,
literature priority nor natural application gain.
